\subsection{Motion, work, and total change}

If \(v(t)\) is velocity, then
\[
\int_a^b v(t)\,dt
\]
is displacement. Distance requires integrating speed \(|v(t)|\), so sign changes must be handled.

\begin{examplebox}[Displacement and distance]
Let \(v(t)=t-1\) on \([0,2]\). Then
\[
\int_0^2 (t-1)\,dt=0,
\]
so the net displacement is zero. But the particle moves one-half unit before \(t=1\) and one-half unit after \(t=1\), so the total distance is
\[
\int_0^2 |t-1|\,dt=1.
\]
\end{examplebox}

If a force \(F(x)\) acts along a line, the work done from \(x=a\) to \(x=b\) is
\[
W=\int_a^b F(x)\,dx.
\]

\begin{examplebox}[Work by a linearly increasing force]
For \(F(x)=2x+1\) from \(x=0\) to \(x=3\),
\[
W=\int_0^3(2x+1)\,dx
=\left[x^2+x\right]_0^3=12.
\]
If force is measured in newtons and distance in metres, the result is \(12\) joules.
\end{examplebox}

\begin{corebox}[Total change from a rate]
Whenever \(Q'(t)\) is a rate of change,
\[
Q(b)-Q(a)=\int_a^b Q'(t)\,dt.
\]
This single principle unifies displacement from velocity, accumulated flow from a flow rate, mass from a linear density, and work from force.
\end{corebox}
